\subsection{Variable jump height}

O této funkci tolik lidí neví, jelikož není na první pohled tak výrazná jako třeba \textit{Wall jump}. To však nic nemění na tom, jak signifikantní rozdíl v uživatelské přívětivosti dělá [4, 5]. Zjednodušeně jde o funkci, která zajišťuje, že když hráč při skoku drží input pro skok delší dobu, tak jeho postava vyskočí výše, než kdyby input jen rychle stiskl. Ve hrách se to často používá právě proto, že jsou někdy v úrovních skoky, kde hráč potřebuje skočit spíše nízko, jinak by například zemřel při kolizi se stropem. Zároveň to podporuje lidskou přirozenost držet input déle při atmosféricky napjatých situacích. Když hráč potřebuje přeskočit delší propast, tak bude delší dobu držet input pro skok.